% SPDX-FileCopyrightText: 2026 Bart van de Lint
% SPDX-License-Identifier: AGPL-3.0-or-later
\section{How sure the comparison is}
\label{sec:certainty}

The savings of \cref{sec:results} come from one model on a sample of the weather. This section moves the router's grid, departures and forecast, each input, each goal, and the few inputs that can change which rig a shipowner should choose.

\subsection{Grid, departures and forecast}
\label{sec:res_robustness}

Each cost saving averages a sample of departures. Against all departures of the three years that leaves an error of 0.4 to 1.0 points (\cref{sec:routes}), so a kite and rotors closer than that are level. \tbd{This draft's stages \SI{720}{nmi} apart lower the cost savings by a further 1.7 points on average for the kite, 1.5 for the rotors and 0.7 for the wing sails against the corridor of the results, whose own error \cref{sec:routing} gives.}

In the base case the router knows the weather of the whole passage, as if its forecast were perfect. No master has that. \tbd{Re-planned daily on the archived forecasts (\cref{fig:forecast}), the wing sails save 5.0\,\% of the cost instead of 7.4\,\% with perfect foresight and the rotors 10.7\,\% instead of 14.4\,\%, keeping 67 and 74\,\% of their cost saving. The kite saves 1.9\,\% of the cost instead of 10.4\,\% and keeps 18\,\%. Perfect foresight therefore favours the kite, and on forecasts it falls behind the wing sails as well as the rotors on this route. It flies 25\,\% of the time at sea instead of 40\,\%, in flights less than half as long. Of its fuel cost saving it keeps 42\,\%, and as its capital and wear stay, its cost saving falls further.}

\tbd{Changing one thing at a time splits the kite's loss, 9.3\,\% of its cost on forecasts. About 5\,\% is the course planned on forecast wind, 1.9\,\% real waves in the plan and 2.2\,\% choosing each step's speed in the weather met, against at most 0.4\,\% for the deck rigs. That choice, blind to the weather ahead, costs the ship without a rig 1.2\,\%, which is why the ship keeps the plan's speed. These parts do not add exactly, and real waves in the plan move the rotors by 2.3\,\%, a measure of the noise of a comparison over 12 departures. Without the bias correction the kite costs 5.1\,\% more again and saves less than no rig. No part of the kite's loss comes from how the model flies it. The kite gains most from a router that finds wind and needs that wind for whole flights, so it loses most when the forecast misplaces it. All these numbers come from a draft run on e-methanol, before this study moved to e-ammonia.} The forecast case is still a bound and not a master's practice: the router plans on one forecast run as if it were certain, without the ensemble's spread or a margin for a kite that may not fly, and re-plans once a day. A router that weighs the kite's flight against the forecast's uncertainty could recover part of its loss.

\begin{figure}[tbp]
\centering
\plot{forecast}
\caption{Cost saving of a \SI{3200}{m^2} kite, \SI{1500}{m^2} of wing sails and \SI{800}{m^2} of rotors on the \SI{140}{m} e-ammonia ship to Guadeloupe, routed with perfect foresight and re-planned every 24 hours on the archived GEFSv12 forecasts, each against the ship without a rig routed the same way (a); and the root-mean-square error against ERA5 along the route of the bias-corrected forecast and of the weather of the same date in another year, by lead time, for the vector wind at 10~m (b) and the significant wave height (c). \tbd{Bars from a draft run on e-methanol.}}
\label{fig:forecast}
\end{figure}

\subsection{Inputs}
\label{sec:res_inputs}

\Cref{fig:sensitivity} moves every input with a range, one at a time, \SI{20}{\%} down and up from its base on the \SI{140}{m} e-ammonia ship to Guadeloupe, and shows the twelve that move any rig most. The same step for every input shows how strongly the result answers to each, whatever is known about it. How uncertain each input is, is its range in \cref{tab:inputs}.

The kite answers most to its own limits and to the wind: the highest wave in which it can still launch moves its cost saving between 10.6 and 19.2\,\%, and the density of the air by almost as much. The wing sails and the rotors answer most to the fuel bill they cut: a less efficient propeller or engine burns more fuel, which raises the rotors' cost saving from 16.2 to 20.2\,\%. \tbd{The figure comes from an earlier version of the model than the class study, which puts the kite at 13.8\,\% here against 15.9\,\% in \cref{tab:inputs}. The final run puts both on one model.}

\Cref{tab:inputs} gives the study at the bounds of each range, on half of the departures, where single inputs move a cost saving by up to 12 points. On the \SI{140}{m} ship the kite and the rotors are level at the base, 15.9 against 16.2\,\%, so any input that helps the kite puts it ahead. On the Kamsarmax the rotors lead, 8.9 against 6.7\,\%. They lose the lead at a high discount rate, a high maintenance share (the yearly maintenance as a share of the rig's capital cost), a high rotor price per area, a low rotor lift, a low kite flight drag or a high wind shear exponent. On the Neo-Panamax container ship the kite leads, 15.9 against 12.3\,\%, and only the lowest bound of the highest waves for a launch puts the rotors ahead.

\begin{figure*}[tbp]
\centering
\plot[\textwidth]{sensitivity}
\caption{Cost saving of a \SI{3200}{m^2} kite (a), \SI{1500}{m^2} of wing sails (b) and \SI{800}{m^2} of rotors (c) on the \SI{140}{m} e-ammonia ship to Guadeloupe with one input \SI{20}{\%} below (blue) or above (orange) its base, for the twelve inputs that move any rig most, in the same order and on the same axis in every panel; the grey line is the saving at the base inputs.}
\label{fig:sensitivity}
\end{figure*}

\subsection{If a goal is missed}
\label{sec:res_missed}

The cost savings above assume that every goal is reached. \Cref{fig:missed_goals} shows, for three ships with their rig sizes held, how each rig's cost saving moves with the canopy life, the canopy price and the share out of service, the share of round trips on which a unit does not work though the wind would allow it. \Cref{tab:missed_goals} gives where each rig stops paying for itself, its cost saving falling to zero, and where it comes level with its rival, the best of the other rigs, if one goal is missed and the others are met. The development that reaching a goal takes is not charged to any rig, so these limits show what development may cost, not whether it pays. \tbd{All numbers in this subsection and its table are draft, from 6 departures on e-methanol where this paper now runs e-ammonia. The final run on 12 replaces them. That run's canopy life goal was \num{4000} flight hours, at which the table and the limits on price, service and crashes hold. The kite's cost savings at the goal of \num{3000} and at half of it are those the run flew at these lives.}

\tbd{At the goals the rotors lead on the \SI{140}{m} ship, 14.4 against 10.6\,\%, and on the Kamsarmax, 6.0 against 3.2\,\%. The kite leads on the Neo-Panamax, 10.4 against 9.3\,\%. The wing sails come last on all three.}

\begin{figure*}[tbp]
\centering
\plot[\textwidth]{missed_goals}
\caption{Cost saving of each rig on three ships against the canopy life, the canopy price and the share of its units' round trips each rig is out of service, the other goals met and the rig sizes held. The kite's saving is swept over the canopy inputs. The dashed lines are the wing sails and rotors at their goals. Vertical lines mark the goals, the dotted line today's canopy price and the shaded band the range of the canopy life input. Where a curve crosses zero the rig stops paying. Where it crosses the next rig's line it falls behind it (\cref{tab:missed_goals}).}
\label{fig:missed_goals}
\end{figure*}

All limits come from one size sweep over the same 12 departures, around the size the class study gives each rig: half to twice for wings and rotors as far as they fit on deck, and a quarter to four times for the kite up to the search bound of \SI{16000}{m^2}. Panel (a) of the table holds each rig at its cheapest size of that sweep and does not resize it as a goal is missed, so a limit says how far a rig already chosen can fall short. Panel (b) gives the crash limits at the same size. \tbd{The cheapest kite lies inside the sizes swept on all three ships. The cheapest wings and rotors are the largest that fit on deck on the \SI{140}{m} ship and the Neo-Panamax, and lie inside the sizes swept on the Kamsarmax.}

\textbf{Canopy life and price.} The top row of \cref{fig:missed_goals} gives the kite's cost saving over canopy lives from 500 to \num{16000} flight hours, so a reader can take the result at any life. \tbd{A canopy of \num{1500} flight hours, half the goal, would cost 1.3 to 2.8 points of cost saving on the three ships, and staying at today's canopy price of USD~234/m$^2$ rather than reaching USD~150/m$^2$ 0.6 to 1.3 points. Within the published range of \numrange{1200}{6900} flight hours (\cref{sec:canopy_life}) canopy life alone reverses kite and rotors only on the Neo-Panamax, where the kite falls behind them below about \num{1700} flight hours. On the other two ships the kite trails the rotors at every life swept. Only on the Kamsarmax does the kite stop paying within the lives and prices swept.} \Cref{sec:check_missed} gives how much of this the routes changing between the two prices explain, and what it is worth to leave the canopy and tether, which wear already pays for, out of maintenance.

\tbd{\textbf{Out of service.} An outage shared by all rigs does not change which leads. Between kite and rotors the order changes only if the leader alone is out of service on 15 to 19\,\% of its round trips, not below the upper bound of 15\,\%, so the order does not rest on the service goal. If all units fail together the rotors fall behind the kite at 12 and 14\,\% on the \SI{140}{m} ship and the Kamsarmax. Every rig keeps paying far beyond any reported outage, until 40 to 97\,\% of round trips, and the wings on the Kamsarmax until 13\,\%.}

\textbf{Crashes.} The goal is no crash, so none is charged. No source gives a crash rate for any rig, so panel (b) gives the crashes a year at which each rig stops paying and at which it comes level with its rival, which needs no rate. A kite that crashes costs its control unit, the steering pod at the end of its tether, and half of its canopy and tether, as wear already charges the other half. A lost wing or rotor costs one new unit with its installation. \tbd{At its cheapest size the kite stops paying at 0.76 to 5.4 crashes a year on the three ships, one per \num{3000} to 560 flight hours, and SkySails reported more than \num{1800} hours of air time over all its kites together \citep{skysails2018}. On the Neo-Panamax it comes level with the rotors sooner, at 0.71 crashes a year. The rotors lose their lead over the kite at 0.27 lost units a year on the \SI{140}{m} ship and 0.23 on the Kamsarmax, one in about four years.} The table gives what each rig can take, not which rig crashes more often.

These limits are upper bounds. They charge money only: the time out after a crash is the share out of service of panel (a), charged separately, and recovery and repair are left out for every rig. Damage costing a tenth of a lost unit is allowed ten times as often. ``Level'' holds the rival free of events. A limit smaller than the error of the cost saving (\cref{sec:routing}) means none. \tbd{Kite and rotors stand further apart than that error on the \SI{140}{m} ship and the Kamsarmax and no further on the Neo-Panamax.}

\begin{table*}[tbp]
\centering
\footnotesize
\caption{\tbd{Draft numbers, from a run at a canopy life goal of \num{4000} flight hours, so the kite's cost saving and the limits against a rival hold at that life.} Where each rig stops paying (pays no more) and comes level with its rival, the best of the other rigs at the goals (level), if one goal is missed and the others are met: the \SI{140}{m} e-ammonia ship to Guadeloupe, the Kamsarmax bulk carrier on VLSFO to Tubar\~ao and the Neo-Panamax container ship on e-ammonia to New York, all on 12 departures. (a) Each rig at its cheapest size, with the points of saving by which it leads its rival at the goals, negative where it trails. For the leading rig a level limit is where it loses the lead, for the others where it would draw level. Out of service counts a unit's round trips. Dashes mark a limit that does not apply, a rig already behind its rival or not paying at the goals, or a limit outside the range swept. Blue shading marks positive cost savings and leads and limits the goal can be missed by, red negative ones and limits the goal falls short of. (b) Crashes a year of each rig at the same size, with its flight hours a year, shaded blue by the room they leave; money only, with no time out, recovery or repair, and the rival has no events.}
\label{tab:missed_goals}
\InputIfFileExists{sections/missed_goals_table}{}{\tbd{Missed-goal table pending the routed run on three ships.}}
\end{table*}

\subsection{Which rig, when}
\label{sec:res_situations}

\Cref{fig:situations} changes one input at a time over its real range on the \SI{140}{m} e-ammonia ship to Guadeloupe and on a Kamsarmax bulk carrier to Tubarão, each time with the best of a few sizes of each rig.

The fuel price decides most. On VLSFO between \SIrange{8}{30}{USD/GJ} no rig pays at the lowest price. Above it the kite leads on the \SI{140}{m} ship, and on the Kamsarmax up to about \tbd{14}\,USD/GJ, with the rotors ahead above it. On e-ammonia from 30 to \SI{120}{USD/GJ} the kite leads on the \SI{140}{m} ship, by \tbd{0.2 to 1.5} points, and the rotors on the Kamsarmax, by \tbd{2.6 to 5.2} points. At \SI{30}{USD/GJ} the two drives give savings about a point apart, so they do not join into one curve.

A higher cargo value raises the interest the cargo costs while at sea and with it the cost of time. The router sails faster and every cost saving shrinks, the kite's fastest, and the rotors pass it at about \tbd{6000}\,USD/t. A higher discount rate weighs more on the rotors and wings, whose cost is mostly capital, than on the kite.

No change to the wing sails alone puts them ahead. At a quarter of their price they save \tbd{13.0}\,\% of the cost on the \SI{140}{m} ship against the kite's \tbd{14.9}\,\%. A higher limit on the rotor surface speed, the speed of the spinning cylinder's skin, moves the cost saving by \tbd{0.1} points or less between \SIrange{35}{65}{m/s}, so in this model the surface speed does not limit the rotor.

\begin{figure*}[tbp]
\centering
\plot[\textwidth]{situations}
\caption{Cost saving of the best size of each rig, from kites of \SIlist{800;1600;3200;6400}{m^2}, wing sails of \SIlist{800;1500;2400}{m^2} and rotors of \SIlist{400;800;1600}{m^2} on the \SI{140}{m} ship to Guadeloupe (solid) and around the class study's best sizes on a Kamsarmax bulk carrier to Tubarão (dashed), with one input changed at a time from its base (dashed grey line; the panel title gives the base value): (a) the fuel price, VLSFO up to \SI{30}{USD/GJ} and e-ammonia from there; (b) the value of the cargo, on the \SI{140}{m} ship only; (c) the discount rate; (d) the life of the wing sails; (e) the unit and area price of the wing sails together; (f) the highest surface speed of the rotors. All but (a) are on e-ammonia.}
\label{fig:situations}
\end{figure*}
